\documentclass[11pt]{article}

% ---------- Packages ----------
\usepackage[a4paper,margin=1in]{geometry}
\usepackage{amsmath,amssymb,amsthm,mathtools}
\usepackage{hyperref}
\usepackage{enumitem}
\usepackage[numbers]{natbib}

% ---------- Theorem environments ----------
\newtheorem{proposition}{Proposition}
\newtheorem{corollary}[proposition]{Corollary}

\theoremstyle{definition}
\newtheorem{definition}[proposition]{Definition}
\newtheorem{remark}[proposition]{Remark}

% ---------- Commands ----------
\newcommand{\Lmu}{\mathcal{L}_{\mathrm{BA},\mu}}
\newcommand{\Tfa}{T_{\mathrm{fa}}}
\newcommand{\N}{\mathbb{N}}

% ---------- Title ----------
\title{Countable Additivity is Not First-Order Axiomatizable}
\author{Patrick S.\ Mahon}
\date{}

\begin{document}
\maketitle

\begin{quote}
\small
\textbf{2020 Mathematics Subject Classification.}
03C20 (primary); 28A60, 60A10 (secondary).

\smallskip
\textbf{Keywords.}
countable additivity, first-order axiomatizability, ultraproduct,
Boolean algebra, finitely additive measure, \L{}o\'s's theorem.
\end{quote}

\begin{abstract}
In the first-order language of Boolean algebras equipped with a normalized finitely
additive charge, no first-order theory characterizes the models in which the charge
is $\sigma$-additive.  Equivalently, countable additivity is not first-order
axiomatizable among finitely additive probability Boolean algebras.  The proof is a
single ultraproduct argument via \L{}o\'{s}'s theorem: Dirac masses on $\mathcal{P}(\N)$
are individually $\sigma$-additive, but the charge induced on the diagonal copy of
$\mathcal{P}(\N)$ in their ultraproduct is the $\{0,1\}$-valued ultrafilter charge,
which is purely finitely additive --- contradicting \L{}o\'{s}'s theorem if any
first-order axiomatization existed.
\end{abstract}

A property of structures is first-order axiomatizable in a language $L$ if there exists
a set of $L$-sentences whose models are exactly the structures satisfying it.
\L{}o\'{s}'s theorem gives the key diagnostic: if a property fails in an ultraproduct
of structures that individually satisfy it, no first-order axiomatization exists.

Countable additivity distinguishes $\sigma$-additive measures from finitely additive
charges.  That it cannot be expressed by a single first-order sentence is clear: the
condition quantifies over countable sequences, unavailable in first-order logic.  Its
non-first-order character is treated as background knowledge in the probability logic
literature~\cite{Hoover1978,Keisler1985,FaginHalpernMegiddo1990}, but the sharper
question --- whether countable additivity admits a first-order \emph{axiomatization}
by any set of sentences, not merely a single formula --- does not appear to have been
settled by an explicit proof.

We settle it negatively.

\medskip

Let $\Lmu$ be the one-sorted first-order language with:
\begin{itemize}
  \item Boolean operations $\wedge, \vee, \neg, 0, 1$;
  \item for each rational $q \in [0,1] \cap \mathbb{Q}$, unary relation symbols
        $M_{\leq q}(x)$ and $M_{\geq q}(x)$, intended to mean
        $\mu(x) \leq q$ and $\mu(x) \geq q$ respectively.
\end{itemize}

A structure for $\Lmu$ is a Boolean algebra $B$ together with a $[0,1]$-valued finitely
additive probability $\mu$, encoded through the truth of the predicates $M_{\leq q}$
and $M_{\geq q}$.  Using rational comparison predicates rather than a function symbol
into $[0,1]$ avoids introducing a second numeric sort.

\begin{definition}[First-order theory of finitely additive probabilities]
Let $\Tfa$ be the $\Lmu$-theory axiomatizing:
\begin{enumerate}[label=(\roman*)]
  \item $B$ is a Boolean algebra;
  \item normalization: $\mu(0) = 0$ and $\mu(1) = 1$;
  \item monotonicity: $x \leq y \Rightarrow \mu(x) \leq \mu(y)$;
  \item finite additivity: for all rationals $r, s$ with $r + s \leq 1$,
        \[
          x \wedge y = 0 \wedge M_{\leq r}(x) \wedge M_{\leq s}(y)
          \;\Rightarrow\; M_{\leq r+s}(x \vee y),
        \]
        and similarly for lower bounds.
\end{enumerate}
\end{definition}

\noindent
The class of Boolean algebras with normalized finitely additive probability is
first-order axiomatizable in $\Lmu$.  Countable additivity, by contrast, requires
quantification over an arbitrary countable sequence $(a_n)_{n \in \N}$ --- not
available in first-order logic --- and so is not a sentence of $\Lmu$.

\medskip

\begin{proposition}[$\sigma$-additivity is not first-order axiomatizable]
\label{prop:not-fo}
There is no first-order $\Lmu$-theory $T \supseteq \Tfa$ such that, for every
$(B, \mu) \models \Tfa$,
\[
  (B, \mu) \models T \quad\Longleftrightarrow\quad \mu \text{ is $\sigma$-additive}.
\]
Equivalently, countable additivity is not first-order axiomatizable among
finitely additive probability Boolean algebras.
\end{proposition}

\begin{proof}
Let $\delta_n$ be the Dirac probability at $n \in \N$ on the $\sigma$-algebra
$\mathcal{P}(\N)$.  Each $(\mathcal{P}(\N), \delta_n)$ is a model of $\Tfa$ in
which $\delta_n$ is $\sigma$-additive.  Suppose for contradiction that a theory
$T$ as described exists; then each $(\mathcal{P}(\N), \delta_n) \models T$.

Let $\mathcal{U}$ be a nonprincipal ultrafilter on $\N$.  By \L{}o\'{s}'s
theorem \cite{Los1955}, the ultraproduct
$\prod_\mathcal{U}(\mathcal{P}(\N),\delta_n)$ also models $T$.  The charge
induced on the diagonal copy of $\mathcal{P}(\N)$ is
\[
  \ell(A) = \lim_\mathcal{U} \delta_n(A) = \lim_\mathcal{U} \mathbf{1}_{n \in A}.
\]
This is the $\{0,1\}$-valued charge corresponding to $\mathcal{U}$: $\ell(A) = 1$
iff $A \in \mathcal{U}$.  Since $\mathcal{U}$ is nonprincipal, $\ell(\{n\}) = 0$
for every $n$, yet $\ell(\N) = 1$.  Thus the sets $A_k = \{k, k+1, \ldots\}$
satisfy $A_0 \supseteq A_1 \supseteq \cdots$ and $\bigcap_k A_k = \varnothing$,
yet $\ell(A_k) = 1$ for all $k$ (since $A_k \in \mathcal{U}$ for all $k$, as
$\mathcal{U}$ is nonprincipal).  Hence $\ell$ is not $\sigma$-additive.

Consequently the charge on the diagonal copy of $\mathcal{P}(\N)$ fails
$\sigma$-additivity, so the ultraproduct cannot satisfy any theory $T$
axiomatizing $\sigma$-additivity, contradicting \L{}o\'{s}'s theorem applied
to the assumption that each $(\mathcal{P}(\N), \delta_n) \models T$.
\end{proof}

\begin{remark}
The ultraproduct $\prod_\mathcal{U}(\mathcal{P}(\N),\delta_n)$ is an
$\Lmu$-structure in its own right; $\ell$ is the charge induced on the diagonal
copy of $\mathcal{P}(\N)$ inside it.  The diagonal copy is isomorphic to
$\mathcal{P}(\N)$ as a Boolean algebra, and $\ell$ is the $\{0,1\}$-valued
ultrafilter charge on it: a purely finitely additive probability that is not
$\sigma$-additive.
\end{remark}

\medskip

The proof exhibits a purely finitely additive charge --- the ultrafilter charge on
$\mathcal{P}(\N)$ --- as a $\mathcal{U}$-ultralimit of $\sigma$-additive
probabilities.  This shows that the class of $\sigma$-additive structures is not
closed under ultraproducts, hence not first-order axiomatizable.  A natural
follow-up question is which purely finitely additive charges on a Boolean algebra
arise as pointwise ultralimits of $\sigma$-additive probabilities on the same
algebra; the answer depends sensitively on the algebra's structural properties and
remains open in the non-$\sigma$-complete non-atomic case.

\bibliographystyle{plain}
\bibliography{../references}

\end{document}
